%==============================================================================
% ANEXA A - Cod Java (Backend)
% Continute REAL din proiectul ETTIway.
% Singura modificare: parola admin inlocuita cu variabila `initialAdminPassword`
% pentru a nu expune secretul in documentatie.
% Comentariile din interiorul lstlisting sunt fara diacritice (cerinta tehnica).
%==============================================================================

\appendix
\chapter{Cod Java (Backend)}
\label{anexa:java}

Această anexă reunește fragmentele de cod Java care implementează componentele
backend descrise în Capitolele~2 și~3. Numerotarea listing-urilor corespunde
trimiterilor făcute în text.

\section*{Punctul de intrare al aplicației}

\begin{lstlisting}[language=Java, caption={Clasa principală a aplicației Spring Boot}, label={lst:a1}]
package com.example.demo;

import org.springframework.boot.SpringApplication;
import org.springframework.boot.autoconfigure.SpringBootApplication;

@SpringBootApplication
public class DemoApplication {

    public static void main(String[] args) {
        SpringApplication.run(DemoApplication.class, args);
    }
}
\end{lstlisting}

\section*{Acces la date prin JPA}

\begin{lstlisting}[language=Java, caption={Interfata de acces la date pentru utilizatori}, label={lst:a2}]
package com.example.demo.repository;

import com.example.demo.entity.User;
import org.springframework.data.jpa.repository.JpaRepository;
import java.util.Optional;

public interface UserRepository extends JpaRepository<User, Long> {
    Optional<User> findByUsername(String username);
}
\end{lstlisting}

\section*{Entități}

\begin{lstlisting}[language=Java, caption={Entitatea User mapata pe tabelul users}, label={lst:a3}]
package com.example.demo.entity;

import jakarta.persistence.*;

@Entity
@Table(name = "users")
public class User {
    @Id
    @GeneratedValue(strategy = GenerationType.IDENTITY)
    private Long id;

    @Column(unique = true, nullable = false)
    private String username;

    @Column(nullable = false)
    private String password;

    private String role; 

    
    public User() {}

    
    public Long getId() { return id; }
    public void setId(Long id) { this.id = id; }
    public String getUsername() { return username; }
    public void setUsername(String username) { this.username = username; }
    public String getPassword() { return password; }
    public void setPassword(String password) { this.password = password; }
    public String getRole() { return role; }
    public void setRole(String role) { this.role = role; }
}
\end{lstlisting}

\begin{lstlisting}[language=Java, caption={Entitatea MapData pentru stocarea harții exterioare și a planurilor indoor}, label={lst:a4}]
package com.example.demo.entity;

import jakarta.persistence.Column;
import jakarta.persistence.Entity;
import jakarta.persistence.Id;
import jakarta.persistence.Table;
import java.time.LocalDateTime;

@Entity
@Table(name = "map_data")
public class MapData {

    @Id
    private Long id;

    @Column(columnDefinition = "TEXT")
    private String graphJson;

    @Column(columnDefinition = "TEXT")
    private String indoorJson;

    private LocalDateTime updatedAt;

    public MapData() {}

    public Long getId() { return id; }
    public void setId(Long id) { this.id = id; }
    public String getGraphJson() { return graphJson; }
    public void setGraphJson(String graphJson) { this.graphJson = graphJson; }
    public String getIndoorJson() { return indoorJson; }
    public void setIndoorJson(String indoorJson) { this.indoorJson = indoorJson; }
    public LocalDateTime getUpdatedAt() { return updatedAt; }
    public void setUpdatedAt(LocalDateTime updatedAt) { this.updatedAt = updatedAt; }
}
\end{lstlisting}

\section*{Securitate: autentificare și autorizare}

\begin{lstlisting}[language=Java, caption={Definirea componentei de criptare a parolelor (BCrypt)}, label={lst:a5}]
@Bean
public PasswordEncoder passwordEncoder() {
    return new BCryptPasswordEncoder();
}
\end{lstlisting}

\begin{lstlisting}[language=Java, caption={Definirea regulilor de autorizare (SecurityConfig)}, label={lst:security}]
@Configuration
public class SecurityConfig {

    @Bean
    public PasswordEncoder passwordEncoder() {
        return new BCryptPasswordEncoder();
    }

    @Bean
    public SecurityFilterChain filterChain(HttpSecurity http) throws Exception {
        http
            .csrf(csrf -> csrf.disable())
            .authorizeHttpRequests(auth -> auth
                .requestMatchers("/login.html", "/auth.html", "/api/auth/**", "/css/**", "/js/**", "/icons/**", "/Icons/**", "/data/**").permitAll()
                .requestMatchers(HttpMethod.GET, "/api/graph/**").authenticated()
                .requestMatchers(HttpMethod.POST, "/api/graph/**").hasRole("ADMIN")
                .requestMatchers(HttpMethod.DELETE, "/api/graph/**").hasRole("ADMIN")
                .requestMatchers(HttpMethod.GET, "/api/indoor/**").authenticated()
                .requestMatchers(HttpMethod.POST, "/api/indoor/**").hasRole("ADMIN")
                .requestMatchers("/api/admin/**").hasRole("ADMIN")
                .anyRequest().authenticated()
            )
            .formLogin(form -> form
                .loginPage("/login.html")
                .loginProcessingUrl("/login")
                .defaultSuccessUrl("/success", true)
                .permitAll()
            )
            .logout(logout -> logout
                .logoutUrl("/logout")
                .logoutSuccessUrl("/login.html")
                .permitAll()
            );
        return http.build();
    }
}
\end{lstlisting}

\begin{lstlisting}[language=Java, caption={Serviciul de incarcare a utilizatorilor pentru autentificare}, label={lst:a7}]
package com.example.demo.config;

import com.example.demo.entity.User;
import com.example.demo.repository.UserRepository;
import org.springframework.beans.factory.annotation.Autowired;
import org.springframework.security.core.authority.SimpleGrantedAuthority;
import org.springframework.security.core.userdetails.*;
import org.springframework.stereotype.Service;

@Service
public class CustomUserDetailsService implements UserDetailsService {

    @Autowired
    private UserRepository userRepository;

    @Override
    public UserDetails loadUserByUsername(String username) throws UsernameNotFoundException {
        User user = userRepository.findByUsername(username)
                .orElseThrow(() -> new UsernameNotFoundException("User not found"));

        return org.springframework.security.core.userdetails.User
                .withUsername(user.getUsername())
                .password(user.getPassword())
                .authorities(new SimpleGrantedAuthority(user.getRole()))
                .build();
    }
}
\end{lstlisting}

\begin{lstlisting}[language=Java, caption={Initializarea automata a contului de administrator}, label={lst:a8}]
@Configuration
public class DataInitializer {

    @Bean
    CommandLineRunner init(UserRepository userRepository, PasswordEncoder passwordEncoder) {
        return args -> {
            if (userRepository.findByUsername("admin").isEmpty()) {
                User admin = new User();
                admin.setUsername("admin");
                admin.setPassword(passwordEncoder.encode(initialAdminPassword));
                admin.setRole("ROLE_ADMIN");
                userRepository.save(admin);
            }
        };
    }
}
\end{lstlisting}

\section*{Persistarea planurilor indoor}

\begin{lstlisting}[language=Java, caption={Repository-ul pentru entitatea MapData}, label={lst:a9}]
package com.example.demo.repository;

import com.example.demo.entity.MapData;
import org.springframework.data.jpa.repository.JpaRepository;
import org.springframework.stereotype.Repository;

@Repository
public interface MapDataRepository extends JpaRepository<MapData, Long> {
}
\end{lstlisting}



\begin{lstlisting}[language=Java, caption={Controller REST pentru salvarea si incarcarea planurilor indoor}, label={lst:a10}]
package com.example.demo.controller;

import com.example.demo.entity.MapData;
import com.example.demo.repository.MapDataRepository;
import org.springframework.beans.factory.annotation.Autowired;
import org.springframework.http.ResponseEntity;
import org.springframework.web.bind.annotation.*;

import java.time.LocalDateTime;
import java.util.Optional;

@RestController
@RequestMapping("/api/indoor")
public class IndoorMapController {

    @Autowired
    private MapDataRepository mapDataRepository;

    @PostMapping("/save")
    public ResponseEntity<String> saveIndoorMap(@RequestBody String indoorJson) {
        MapData mapData = mapDataRepository.findById(1L).orElse(new MapData());
        mapData.setId(1L);
        mapData.setIndoorJson(indoorJson);
        mapData.setUpdatedAt(LocalDateTime.now());
        mapDataRepository.save(mapData);
        return ResponseEntity.ok("Indoor map saved successfully.");
    }

    @GetMapping("/load")
    public ResponseEntity<String> loadIndoorMap() {
        Optional<MapData> mapDataOpt = mapDataRepository.findById(1L);
        if (mapDataOpt.isPresent() && mapDataOpt.get().getIndoorJson() != null) {
            return ResponseEntity.ok(mapDataOpt.get().getIndoorJson());
        }
        return ResponseEntity.ok("{}");
    }
}
\end{lstlisting}

\begin{lstlisting}[language=Java, caption={Controller REST pentru salvarea, incarcarea si stergerea grafului exterior}, label={lst:a11}]
@RestController
@RequestMapping("/api/graph")
public class GraphController {

    @Autowired
    private MapDataRepository mapDataRepository;

    @PostMapping("/save")
    public ResponseEntity<String> saveGraph(@RequestBody String graphJson) {
        MapData mapData = mapDataRepository.findById(1L).orElse(new MapData());
        mapData.setId(1L);
        mapData.setGraphJson(graphJson);
        mapData.setUpdatedAt(LocalDateTime.now());
        mapDataRepository.save(mapData);
        return ResponseEntity.ok("Graph saved successfully.");
    }

    @GetMapping("/load")
    public ResponseEntity<String> loadGraph() {
        Optional<MapData> mapDataOpt = mapDataRepository.findById(1L);
        if (mapDataOpt.isPresent() && mapDataOpt.get().getGraphJson() != null) {
            return ResponseEntity.ok(mapDataOpt.get().getGraphJson());
        }
        return ResponseEntity.ok("{}");
    }

    @DeleteMapping("/erase")
    public ResponseEntity<String> eraseGraph() {
        if (mapDataRepository.existsById(1L)) {
            mapDataRepository.deleteById(1L);
            return ResponseEntity.ok("Graph erased successfully.");
        }
        return ResponseEntity.ok("No graph found to erase.");
    }
}
\end{lstlisting}